%% ============================================================
%% §4 — Multi-Platform Validation Design
%% RL Control Agent section
%% Describes the three replication/extension experiments and
%% the algorithms used
%% ============================================================

\section{Multi-Platform Validation Design}
\label{sec:validation}\label{sec:control_agent}

The central empirical claim of this paper is that SiliQun is
\emph{neutral and fair} across silicon spin qubit technologies: the
simulator does not artificially favour any one hardware platform, and any
machine-learning algorithm that works on one profile can be applied to any
other without modification.
To test this claim rigorously, we adopt a \emph{replication-based
validation} strategy: we select two independently published
machine-learning experiments on silicon spin qubits, reproduce them
verbatim on SiliQun using the matching hardware profile, and compare our
results to the published values.
We then extend the benchmark to GAA qubits, for which no prior
machine-learning experiment exists, to demonstrate that SiliQun enables
experiments that cannot yet be conducted on physical hardware.

%% ----------------------------------------------------------
\subsection{Replication 1: DQL and SARSA on SiMOS}
\label{subsec:rep1_design}
%% ----------------------------------------------------------

Daraeizadeh~et~al.~\cite{Daraeizadeh2020} applied deep Q-learning (DQL)
and SARSA to the problem of CZ gate synthesis in a two-qubit SiMOS quantum
dot device.
Their Hamiltonian is the exchange-coupled two-qubit ZZ model:
\begin{equation}
  H(t) = \frac{J(t)}{4} \sigma_z^{(1)} \otimes \sigma_z^{(2)}
         + \frac{\varepsilon_0(t)}{2} \sigma_z^{(1)}
         + \frac{\varepsilon_1(t)}{2} \sigma_z^{(2)},
  \label{eq:zz_hamiltonian}
\end{equation}
where $J(t)$ is the exchange coupling and $\varepsilon_0(t)$,
$\varepsilon_1(t)$ are the single-qubit detuning parameters, all
controllable by the RL agent.
The target gate is CZ, and the fidelity metric is the phase-invariant
process fidelity:
\begin{equation}
  F = \frac{1}{d^2} \left| \mathrm{tr}(U_\text{target}^\dagger \,
      U_\text{achieved}) \right|^2,
  \label{eq:phase_fidelity}
\end{equation}
where $d = 4$ for a two-qubit system.
The published result is $F > 0.999$ for both DQL and SARSA.

We replicate this experiment on SiliQun's \texttt{simos\_nominal} profile
with the following faithful reproduction of the original setup:
a discrete action space of 27 parameter combinations
($\varepsilon_0, \varepsilon_1, J \in \{-\Delta, 0, +\Delta\}$ with
$\Delta = 0.1$), an episode length of 100 steps, an $\varepsilon$-greedy
exploration schedule decaying from 1.0 to 0.01 over 5,000 episodes, a
replay buffer of 10,000 transitions, and a target network updated every
100 steps.
The DQL agent uses a two-layer network with 16 hidden units per layer,
consistent with the small network described in the original paper.
Five independent seeds are run for each algorithm.

%% ----------------------------------------------------------
\subsection{Replication 2: DQN for Universal State Preparation on Donor
  Qubits}
\label{subsec:rep2_design}
%% ----------------------------------------------------------

He~et~al.~\cite{He2021} applied a deep Q-network (DQN) to the universal
state preparation (USP) problem in a singlet-triplet ($S$-$T_0$) quantum
dot device.
The $S$-$T_0$ Hamiltonian is:
\begin{equation}
  H(t) = J(t) \sigma_z + h \sigma_x,
  \label{eq:st0_hamiltonian}
\end{equation}
where $J(t)$ is the exchange coupling (controllable) and $h$ is the
hyperfine gradient (fixed at $h = A_\text{hf}/2$).
The task is to prepare an arbitrary target state $|\psi_\text{target}\rangle$
on the Bloch sphere from the initial state $|0\rangle$, with fidelity
$F = |\langle \psi_\text{target} | \psi \rangle|^2 \geq 0.99$.
The published result is $F > 0.99$ for target states on the six cardinal
axes of the Bloch sphere ($|0\rangle$, $|1\rangle$, $|\pm\rangle$,
$|\pm i\rangle$).

We replicate this experiment on SiliQun's \texttt{donor\_electron} profile,
which implements the $S$-$T_0$ physics with the hyperfine coupling
$A_\text{hf} = 117.5$\,MHz drawn from the M\k{a}dzik~et~al.\
characterisation~\cite{Madzik2022}.
The DQN agent uses a two-layer network with 32 hidden units, an episode
length of 50 steps, an $\varepsilon$-greedy schedule decaying from 1.0 to
0.01 over 2,000 episodes, and a replay buffer of 5,000 transitions.
Evaluation is performed on the six cardinal Bloch sphere states following
the He~et~al.\ protocol.
Five independent seeds are run.

%% ----------------------------------------------------------
\subsection{Extension 3: DQN for Universal State Preparation on GAA Qubits}
\label{subsec:ext3_design}
%% ----------------------------------------------------------

No prior machine-learning experiment has been reported for GAA silicon spin
qubits.
We apply the identical DQN USP protocol from Replication~2 to
SiliQun's \texttt{gaa\_nominal} profile, changing only the hardware profile
constructor argument.
The GAA profile has a shorter coherence time ($T_2 = 0.1$\,ms versus
$T_2 = 1$\,ms for donor) and stronger charge noise
($\sigma_\varepsilon = 2.0$\,$\mu$eV versus $0.3$\,$\mu$eV), which we
expect to reduce the achievable fidelity relative to the donor profile.
The degree of this reduction is a direct empirical measure of the
decoherence penalty associated with GAA's CMOS-compatible fabrication
process.
This extension constitutes the first machine-learning experiment on GAA
silicon spin qubits reported in the literature.

%% ----------------------------------------------------------
\subsection{Algorithm Implementations}
\label{subsec:algorithms}
%% ----------------------------------------------------------

All three experiments use standard, well-established RL algorithms
without modification.
The DQL and SARSA implementations follow the original
Daraeizadeh~et~al.\ description: a two-layer fully connected network with
ReLU activations, mean-squared error loss, and the Adam optimiser with
learning rate $\alpha = 10^{-3}$.
SARSA uses on-policy updates (the current policy's action is used for the
TD target), while DQL uses off-policy updates (the greedy action is used
for the TD target).
The DQN implementation follows the He~et~al.\ description with the same
network architecture, experience replay, and target network.

All experiments are run on an NVIDIA RTX~2070 Max-Q GPU (8\,GB VRAM)
using PyTorch~2.6.0 with CUDA~12.4.
The analytical ZZ Hamiltonian exponential is computed in closed form to
achieve the throughput required for 10,000-episode training runs.
All code and hardware profile definitions are released in the
\texttt{siliqun-provider} repository under the MIT licence.

